% GENERATED FILE. Edit canonical lessons, then run npm run generate:books.
% canonical-source-hash: fnv1a64:846dc32bf9fe1174
% canonical-lessons: SA-C291-dhurtah, SA-C291-vancayati, SA-C291-utkocah, SA-C291-kapatam, SA-C291-krtrimah

\chapter{Cheats}
\label{ch:sa-cheats}
\hlchaptermodality{291}

\begin{chapteropening}
\textbf{By the end of this chapter:} \emph{I can say how someone was cheated and what to do if it happens.}

Complete the last lesson of chapter 291: I can say how someone was cheated and what to do if it happens.
\end{chapteropening}

\section[\texorpdfstring{\sk{धूर्तः} (dhūrtaḥ)}{धूर्तः (dhūrtaḥ)}]{\sk{धूर्तः} (dhūrtaḥ) — a rogue; a cheat}

\label{lesson:SA-C291-dhurtah}



\begin{quote}
Before the new one: say the Sanskrit for a piece of jewellery, then the Sanskrit for by force.
\end{quote}

\subsection*{You'll want to know: \sk{धूर्तः}}
\textbf{\sk{धूर्तः}} — \emph{dhūrtaḥ} — \textquotedblleft{}a rogue; a cheat\textquotedblright{}.

\begin{quote}
\sk{सः} \sk{नरः} \sk{धूर्तः} \sk{अस्ति।}
\end{quote}

\begin{quote}
\emph{saḥ naraḥ dhūrtaḥ asti.}
\end{quote}

\begin{quote}
That man is a rogue.
\end{quote}

A masculine noun.

\begin{cousinweb}
\textbf{\sk{धूर्तः}} (\emph{dhūrtaḥ}) comes from the root √dhūrv, \textquotedblleft{}to hurt\textquotedblright{} (MW also names √dhvṛ); MW gives \emph{dhūrta} as cunning and, as a noun, a rogue or swindler (MW).
\end{cousinweb}

\begin{grammarlens}[title={What you've built}]
One more word for telling what happened.
\end{grammarlens}

\subsection*{Your turn}
\begin{itemize}
\raggedright
  \item \emph{Say it:} \emph{dhūrtaḥ}
  \item \emph{Say it:} \emph{dhūrtaḥ}, once more
  \item \emph{Say it:} \emph{saḥ naraḥ dhūrtaḥ asti.} — That man is a rogue.
  \item \emph{Recall:} say the Sanskrit for a piece of jewellery, then the Sanskrit for by force, then say \emph{dhūrtaḥ} again
\end{itemize}

\begin{tcolorbox}[breakable,colback=teal!4,colframe=teal!35!black,title={Before you move on}]
What does \sk{धूर्तः} mean? (\textquotedblleft{}A rogue; a cheat\textquotedblright{}.) Say it once more.
\end{tcolorbox}
\section[\texorpdfstring{\sk{वञ्चयति} (vañcayati)}{वञ्चयति (vañcayati)}]{\sk{वञ्चयति} (vañcayati) — cheats; deceives}

\label{lesson:SA-C291-vancayati}



\begin{quote}
Before the new one: say the Sanskrit for by force, then the Sanskrit for a rogue.
\end{quote}

\subsection*{You'll want to know: \sk{वञ्चयति}}
\textbf{\sk{वञ्चयति}} — \emph{vañcayati} — \textquotedblleft{}cheats; deceives\textquotedblright{}.

\begin{quote}
\sk{धूर्तः} \sk{वृद्धं} \sk{वञ्चयति।}
\end{quote}

\begin{quote}
\emph{dhūrtaḥ vṛddhaṃ vañcayati.}
\end{quote}

\begin{quote}
The rogue cheats the old man.
\end{quote}

Past (he …): \sk{सः} \sk{वृद्धं} \sk{वञ्चितवान्}, \emph{saḥ vṛddhaṃ vañcitavān}.

\begin{cousinweb}
\textbf{\sk{वञ्चयति}} (\emph{vañcayati}) comes from the root √vañc, \textquotedblleft{}to go crookedly\textquotedblright{}; MW gives its causative \emph{vañcayati} as to avoid, to lead astray and to cheat (MW).
\end{cousinweb}

\begin{grammarlens}[title={What you've built}]
One more word for telling what happened.
\end{grammarlens}

\subsection*{Your turn}
\begin{itemize}
\raggedright
  \item \emph{Say it:} \emph{vañcayati}
  \item \emph{Say it:} \emph{vañcayati}, once more
  \item \emph{Say it:} \emph{dhūrtaḥ vṛddhaṃ vañcayati.} — The rogue cheats the old man.
  \item \emph{Recall:} say the Sanskrit for by force, then the Sanskrit for a rogue, then say \emph{vañcayati} again
\end{itemize}

\begin{tcolorbox}[breakable,colback=teal!4,colframe=teal!35!black,title={Before you move on}]
What does \sk{वञ्चयति} mean? (\textquotedblleft{}Cheats; deceives\textquotedblright{}.) Say it once more.
\end{tcolorbox}
\section[\texorpdfstring{\sk{उत्कोचः} (utkocaḥ)}{उत्कोचः (utkocaḥ)}]{\sk{उत्कोचः} (utkocaḥ) — a bribe}

\label{lesson:SA-C291-utkocah}



\begin{quote}
Before the new one: say the Sanskrit for a rogue, then the Sanskrit for cheats.
\end{quote}

\subsection*{You'll want to know: \sk{उत्कोचः}}
\textbf{\sk{उत्कोचः}} — \emph{utkocaḥ} — \textquotedblleft{}a bribe\textquotedblright{}.

\begin{quote}
\sk{यदि} \sk{सः} \sk{उत्कोचं} \sk{याचते}, \sk{तर्हि} \sk{धूर्तः।}
\end{quote}

\begin{quote}
\emph{yadi saḥ utkocaṃ yācate, tarhi dhūrtaḥ.}
\end{quote}

\begin{quote}
If he asks for a bribe, then he is a cheat.
\end{quote}

A masculine noun.

\begin{cousinweb}
\textbf{\sk{उत्कोचः}} (\emph{utkocaḥ}) is built from \emph{ut} and the root √kuc, \textquotedblleft{}to bend, contract\textquotedblright{}; MW's senses run from unwinding to bribery and corruption (MW).
\end{cousinweb}

\begin{grammarlens}[title={What you've built}]
One more word for telling what happened.
\end{grammarlens}

\subsection*{Your turn}
\begin{itemize}
\raggedright
  \item \emph{Say it:} \emph{utkocaḥ}
  \item \emph{Say it:} \emph{utkocaḥ}, once more
  \item \emph{Say it:} \emph{yadi saḥ utkocaṃ yācate, tarhi dhūrtaḥ.} — If he asks for a bribe, then he is a cheat.
  \item \emph{Recall:} say the Sanskrit for a rogue, then the Sanskrit for cheats, then say \emph{utkocaḥ} again
\end{itemize}

\begin{tcolorbox}[breakable,colback=teal!4,colframe=teal!35!black,title={Before you move on}]
What does \sk{उत्कोचः} mean? (\textquotedblleft{}A bribe\textquotedblright{}.) Say it once more.
\end{tcolorbox}
\section[\texorpdfstring{\sk{कपटम्} (kapaṭam)}{कपटम् (kapaṭam)}]{\sk{कपटम्} (kapaṭam) — fraud; a trick}

\label{lesson:SA-C291-kapatam}



\begin{quote}
Before the new one: say the Sanskrit for cheats, then the Sanskrit for a bribe.
\end{quote}

\subsection*{You'll want to know: \sk{कपटम्}}
\textbf{\sk{कपटम्}} — \emph{kapaṭam} — \textquotedblleft{}fraud; a trick\textquotedblright{}.

\begin{quote}
\sk{एतत्} \sk{सर्वं} \sk{कपटम्} \sk{अस्ति।}
\end{quote}

\begin{quote}
\emph{etat sarvaṃ kapaṭam asti.}
\end{quote}

\begin{quote}
All this is a fraud.
\end{quote}

A neuter noun.

\begin{cousinweb}
\textbf{\sk{कपटम्}} (\emph{kapaṭam}) MW gives \emph{kapaṭa} as fraud, deceit and cheating, and reports a commentator's derivation from √kamp, \textquotedblleft{}to tremble\textquotedblright{}, so its origin is uncertain (MW).
\end{cousinweb}

\begin{grammarlens}[title={What you've built}]
One more word for telling what happened.
\end{grammarlens}

\subsection*{Your turn}
\begin{itemize}
\raggedright
  \item \emph{Say it:} \emph{kapaṭam}
  \item \emph{Say it:} \emph{kapaṭam}, once more
  \item \emph{Say it:} \emph{etat sarvaṃ kapaṭam asti.} — All this is a fraud.
  \item \emph{Recall:} say the Sanskrit for cheats, then the Sanskrit for a bribe, then say \emph{kapaṭam} again
\end{itemize}

\begin{tcolorbox}[breakable,colback=teal!4,colframe=teal!35!black,title={Before you move on}]
What does \sk{कपटम्} mean? (\textquotedblleft{}Fraud; a trick\textquotedblright{}.) Say it once more.
\end{tcolorbox}
\section[\texorpdfstring{\sk{कृत्रिमः} (kṛtrimaḥ)}{कृत्रिमः (kṛtrimaḥ)}]{\sk{कृत्रिमः} (kṛtrimaḥ) — fake; artificial}

\label{lesson:SA-C291-krtrimah}



\begin{quote}
Before the new one: say the Sanskrit for a bribe, then the Sanskrit for fraud.
\end{quote}

\subsection*{You'll want to know: \sk{कृत्रिमः}}
\textbf{\sk{कृत्रिमः}} — \emph{kṛtrimaḥ} — \textquotedblleft{}fake; artificial\textquotedblright{}.

\begin{quote}
\sk{एतत्} \sk{आभरणं} \sk{कृत्रिमम्} \sk{अस्ति।}
\end{quote}

\begin{quote}
\emph{etat ābharaṇaṃ kṛtrimam asti.}
\end{quote}

\begin{quote}
This jewellery is fake.
\end{quote}

\begin{cousinweb}
\textbf{\sk{कृत्रिमः}} (\emph{kṛtrimaḥ}) comes from the root √kṛ, \textquotedblleft{}to make\textquotedblright{}: MW gives \emph{kṛtrima} as made by art, artificial, and also falsified (MW).
\end{cousinweb}

\begin{grammarlens}[title={What you've built}]
One more word for telling what happened.
\end{grammarlens}

\subsection*{Your turn}
\begin{itemize}
\raggedright
  \item \emph{Say it:} \emph{kṛtrimaḥ}
  \item \emph{Say it:} \emph{kṛtrimaḥ}, once more
  \item \emph{Say it:} \emph{etat ābharaṇaṃ kṛtrimam asti.} — This jewellery is fake.
  \item \emph{Recall:} say the Sanskrit for a bribe, then the Sanskrit for fraud, then say \emph{kṛtrimaḥ} again
\end{itemize}

\begin{tcolorbox}[breakable,colback=teal!4,colframe=teal!35!black,title={Before you move on}]
What does \sk{कृत्रिमः} mean? (\textquotedblleft{}Fake; artificial\textquotedblright{}.) Say it once more.
\end{tcolorbox}
